\documentclass[10pt,a4paper]{amsart}
\usepackage[margin=19mm]{geometry}
\usepackage{amsmath,amssymb,mathtools}
\usepackage[scr=rsfs,cal=euler]{mathalfa}
\usepackage{xfrac}
\usepackage[unicode,hidelinks,bookmarks=false]{hyperref}
\newcommand{\ie}{i.e.\ }
\newcommand{\cf}{cf.\ }
\newcommand{\resp}{respectively\ }
\newcommand{\Subsection}{\subsection}
\providecommand{\nobreakdash}{\nobreak\hbox{-}}
\setlength{\emergencystretch}{2em}
\multlinegap0pt
\usepackage{bm}
\usepackage{bbm}
\usepackage{dsfont}
\usepackage{xspace}
\usepackage{cancel}
\usepackage{mathtools}
\usepackage{cleveref}
\usepackage{amsthm}
\makeatletter
\@ifundefined{theorem}{\newtheorem{theorem}{Theorem}}{}
\@ifundefined{lemma}{\newtheorem{lemma}[theorem]{Lemma}}{}
\@ifundefined{proposition}{\newtheorem{proposition}[theorem]{Proposition}}{}
\makeatother
\title[On Diophantine properties for values of Dedekind zeta functi]{On Diophantine properties for values of Dedekind zeta functions}
\author{Jerson Caro, Fabien Pazuki, Riccardo Pengo}
\date{Source: arXiv:2502.20910v2 (CC BY 4.0)}
\begin{document}
\maketitle

\noindent\emph{Source and licence.} On Diophantine properties for values of Dedekind zeta functions. Version arXiv:2502.20910v2 (CC BY 4.0), distributed under the Creative Commons Attribution 4.0 International licence, as recorded in the arXiv licence metadata for this exact version and shown on the abstract page https://arxiv.org/abs/2502.20910v2. Full rights notice: licenses/arxiv-2502.20910-CC-BY-4.0.txt.\par\smallskip

\noindent\textit{Editorial scope.} Counted unit: the Theorem whose source label is \texttt{M1alpha} in the article (source0.tex lines 1397--1404, source label \texttt{M1alpha}), delivered together with the article's own complete original proof of it (source0.tex lines 1405--1420, one \texttt{proof} environment of 16 lines). No mathematics is written, repaired or re-derived: statement and proof are the article's own text line for line. Required context retained uncounted: eq:FirstMoment\_right (lines 1213--1216); L-function approximation (lines 1266--1276); M1 approximation (lines 1347--1353). Every definition, display and label of those lines is retained unchanged. Omitted from this package and not redistributed: 1--1212, 1217--1265, 1277--1346, 1354--1396, 1421--1675. Nothing inside a retained excerpt is omitted or altered: every retained body is delivered exactly as the article's lines are written, so no omission receipt of any kind accompanies this package. Citations: the retained excerpts carry no citation command at all, so no bibliography entry of the article is transcribed and no reference is redistributed. Reference adaptation: none was needed and none was made; every reference inside the delivered unit resolves inside it, so no unresolved reference or citation is produced. Printed slips kept literal: no printed word, symbol, hypothesis or number of the counted statement or of its proof was altered. Layout and class: the article's own class is \texttt{amsart} with packages including inputenc, fontenc, mathpazo, stackengine, amsmath, amsthm, amssymb, mathtools, soul, wrapfig, mathdots, graphicx, tikz, tikz-cd; the wrapper is the lane's pinned \texttt{amsart} editorial wrapper at 10pt on A4, which restates the theorem-like environments the retained text needs and adds the article's own macro definitions that the retained text needs (none), with extra wrapper packages (bm, bbm, dsfont, xspace, cancel, mathtools, cleveref). The delivered numbering is the unit's own. Assets: no retained span contains a figure environment, a graphics-inclusion command, a PDF-page inclusion, a table or a TikZ picture; no image file is copied into this package, the package declares no asset dependency and no image file is redistributed with it. Licence: version arXiv:2502.20910v2 of this article is distributed under the Creative Commons Attribution 4.0 International licence, as recorded in the arXiv licence metadata for this exact version and shown on the abstract page https://arxiv.org/abs/2502.20910v2. A full rights notice is delivered at licenses/arxiv-2502.20910-CC-BY-4.0.txt. Characters: every retained excerpt is pure ASCII, so the delivered engine, XeLaTeX with its default Latin Modern text font, reports no missing character.
\begin{align}
    M_1(R,X,\sigma) &= \sum_{\frac{X}{8}\leq d\leq \frac{5X}{16}}\mu(2d)^2R(8d)^2 L(\sigma,\chi_{8d}) \label{eq:FirstMoment_right} \\
    M_2(R,X,\sigma) &= \sum_{\frac{X}{8}\leq d\leq \frac{5X}{16}} \mu(2d)^2 R(8d)^2 (L(\sigma,\chi_{8d}))^2. \label{eq:SecondMoment_right} 
\end{align}

\begin{proposition}
\label{L-function approximation}
For every $\sigma \in (\frac{1}{2},1)$ and $\delta \in \left(0, \sigma-\frac{1}{2}\right)$ there exists two positive constants $c_8(\sigma,\delta)$ and $c_9(\delta)$ such that for every $X \geq c_9(\delta)$ we have that
\begin{equation}\label{error}
\sum_{\frac{X}{8}\leq d\leq \frac{5X}{16}}\mu(2d)^2R(8d)^2\left|L(\sigma,\chi_{8d})-\sum_{n=1}^{\infty}\frac{\chi_{8d}(n)}{n^\sigma} e^{-\frac{n}{X^2}}\right|^2 \leq c_8(\sigma,\delta) X^{1-4\delta}\prod_{p\text{ prime}}\left(1+f(p)^2+4\frac{f(p)}{p^{\sigma-\delta}}\right).
\end{equation}
In particular, we have that
\[
\left\lvert M_1(R,X,\sigma) - \sum_{\frac{X}{16}\leq d\leq \frac{X}{8}}\mu(2d)^2R(8d)^2\sum_{n=1}^{\infty}\frac{\chi_{8d}(n)}{n^\sigma} e^{-\frac{n}{X^2}} \right\rvert \leq c_8(\sigma,\delta) X^{1-4\delta}\prod_{p\text{ prime}}\left(1+f(p)^2+4\frac{f(p)}{p^{\sigma-\delta}}\right).
\].
\end{proposition}

\begin{lemma}\label{M1 approximation}
Let $\sigma \in \left( \frac{1}{2},1 \right)$. Then, we have that
\begin{equation}\label{eq M1 approx}
\sum_{\frac{X}{8}\leq d\leq \frac{5X}{16}}\mu(2d)^2R(8d)^2\sum_{n=1}^{\infty}\frac{\chi_{8d}(n)}{n^\sigma} e^{-\frac{n}{X^2}} \sim c_{13}(\sigma) X\prod_{p\text{ prime}}\left(1+f(p)^2-2\frac{f(p)}{p^{\sigma}}\right)   
\end{equation}
as $X \to +\infty$, where $c_{13}(\sigma)$ is some positive constant depending only on $\sigma$.
\end{lemma}

\begin{theorem}\label{M1alpha}
Let $\sigma \in \left( \frac{1}{2},1 \right)$. Then, there exist two positive constants $c_{14}(\sigma)$ and $c_{15}(\sigma)$ such that
\begin{equation*}
M_1(R,X,\sigma)
\geq c_{14}(\sigma) X\prod_{p\text{ prime}}\left(1+f(p)^2-2\frac{f(p)}{p^{\sigma}}\right),
\end{equation*}
for every $X \geq c_{15}(\sigma)$, where $M_1(R,X,\sigma)$ is defined as in \eqref{eq:FirstMoment_right}.
\end{theorem}

\begin{proof}
By \cref{L-function approximation} and \cref{M1 approximation} it is enough to prove that
\begin{equation} \label{eq:right_firstmoment_final}
    \frac{1}{X^{4\delta}}\prod_{p\text{ prime}}\frac{\left(1+f(p)^2+4\frac{f(p)}{p^{\sigma-\delta}}\right)}{\left(1+f(p)^2-2\frac{f(p)}{p^{\sigma}}\right)}=o(1),
\end{equation}
for some positive real number $\delta$ satisfying $\delta<\sigma-\frac{1}{2}$.
Notice that 
\begin{equation}\label{asymM1}
\begin{aligned}
\prod_{p \ \text{prime}} \frac{\left(1+f(p)^2+4\frac{f(p)}{p^{\sigma-\delta}}\right)}{\left(1+f(p)^2-2\frac{f(p)}{p^{\sigma}}\right)}
&=\prod_{p \ \text{prime}}\left(1+\frac{4\frac{f(p)}{p^{\sigma-\delta}}+2\frac{f(p)}{p^\sigma}}{1+f(p)^2-2\frac{f(p)}{p^{\sigma}}}\right) \\
&\ll \exp\left((6+o(1))L\sum_{\substack{p \ \text{prime} \\ p\geq L^{\frac{1}{\sigma}}}}\frac{1}{p^{2\sigma-\delta}}\right) \\ &= \exp\left((6+o(1))\frac{L}{ L^{\frac{2\sigma-\delta-1}{\sigma}}\log( L^{\frac{1}{\sigma}})}\right)\\ &=\exp\left(\left(12\sigma+o(1)\right)\frac{\log^{\frac{1+\delta-\sigma}{2\sigma}} X}{\log\log X}\right).
\end{aligned}
\end{equation}
Notice that $1+\delta-\sigma<1$ and therefore $\exp\left(\left(12\sigma+o(1)\right)\frac{\log^{\frac{1+\delta-\sigma}{2\sigma}} X}{\log\log X}\right) = o(X^{\varepsilon})$ for every $\varepsilon>0$, and in particular for $\varepsilon=4\delta$, this computation shows that \eqref{eq:right_firstmoment_final} holds true, and therefore concludes the proof.
\end{proof}

\end{document}
